\documentclass[11pt]{article}
\usepackage[margin=1in]{geometry}
\usepackage{amsmath,amssymb,amsthm}
\usepackage{booktabs}
\usepackage{xurl}
\usepackage[hidelinks]{hyperref}
\newtheorem{theorem}{Theorem}
\newtheorem{lemma}[theorem]{Lemma}
\newtheorem{proposition}[theorem]{Proposition}
\newcommand{\OPT}{\mathrm{OPT}}
\newcommand{\h}{\tfrac12}
\setlength{\emergencystretch}{3em}

\title{Optimal one-round agreement through stochastic channels for four and five processes}
\author{Yitzchak Shmalo$^{\,a}$\thanks{Email address:
\texttt{yitzchak.shmalo@gmail.com} (Yitzchak Shmalo).}\\
{\small $^{a}$Einstein Institute of Mathematics, The Hebrew University of
Jerusalem, Givat Ram,}\\
{\small Jerusalem, Israel}}
\date{November 2026}

\begin{document}
\maketitle

\begin{abstract}
In the agreement optimization problem of Fraigniaud, Patt-Shamir and Rajsbaum, processes with binary inputs exchange one round of messages, each delivered independently with probability $p$, and then decide real values; when all inputs agree, every process must decide the common input. The cost of an algorithm is its expected discrepancy, the largest difference between two decisions, on the worst input assignment. The optimal one-round cost was known for two and three processes. We determine it for four and five processes and every $p$. For four processes the optimum is given by four formulas on four intervals of $p$, and for five processes by seven, with algebraic endpoints; on the upper intervals the optimal decisions depend on $p$. At $p=1/2$ the optima are 941/1936 and 394955/851328. The lower bounds are proved in exact arithmetic by weighting two input assignments and checking every decision vector whose entries are 0, 1/2 or 1.
\end{abstract}

\noindent\textit{Keywords:} approximate agreement, distributed computing, message loss, minimax, piecewise linear optimization.

\noindent\textit{2020 Mathematics Subject Classification:} 68W15, 68Q87, 90C47.

\section{Introduction}\label{sec:intro}

Fraigniaud, Patt-Shamir and Rajsbaum \cite{FPR} introduced a stochastic version of approximate agreement \cite{DLPSW}. There are $n$ synchronous processes on a complete network. Process $i$ holds an input $x_i\in\{0,1\}$. In each round it sends a message to every other process, and each message is delivered independently with probability $p$; in a one-round algorithm a process may be taken to send its input. After a fixed number of rounds every process $i$ decides a value $d_i\in[0,1]$. Validity requires that if all inputs equal $x$ then every process decides $x$. The discrepancy of an execution is $\max_i d_i-\min_i d_i$. The cost of a deterministic algorithm is its expected discrepancy, maximized over input assignments. We write $q=1-p$ and $\OPT_n(p)$ for the least cost of a one-round algorithm for $n$ processes.

For two processes Fraigniaud, Patt-Shamir and Rajsbaum proved that $\OPT_2(p)=\min\{q,\,p^2+q^2\}$ \cite[Theorem~3]{FPR}. The optimum is attained for $p\le\h$ by their algorithm Flip Value (FV), in which a process that receives the other value decides it, and for $p\ge\h$ by their Agreed Meeting Point AMP$(a)$, in which such a process decides $a$, for any $a\in[0,1]$. For two processes and any number of rounds some optimal algorithm decides only the values $0$ and $1$ \cite[Section~4.2]{FPR}. This fails for three processes. There they write the one-round problem as a linear program in the decisions, with one constraint for each input assignment, and solve it \cite[Section~5.1]{FPR}. The optimum $\OPT_3$ has two thresholds, $\theta\approx0.3461$ and $\h$, and on each of the three intervals one algorithm, the same for every $p$ in the interval, is optimal \cite[Theorem~13]{FPR}. For large $n$ the expected discrepancy of AMP$(a)$ after one round tends to $0$ on the input assignments in which each of the two values is held by $\omega(\log_{1/q}n)$ processes \cite[Theorem~14]{FPR}, and a two-round algorithm achieves this on every input assignment \cite[Theorem~15]{FPR}. They ask for the optimal algorithms for every $n>2$ \cite[Section~7]{FPR}. We determine $\OPT_4$ and $\OPT_5$ for every $p$ (Theorems~\ref{thm:four} and~\ref{thm:five}).

The cases $n=4$ and $n=5$ differ from $n\le3$ in one respect. Up to a permutation of the processes an input assignment is determined by its number $k$ of ones, and it suffices to consider algorithms given by a vector $y\in[0,1]^m$ of decisions (Lemma~\ref{lem:sym}), with $m=6$ for $n=4$ and $m=10$ for $n=5$. Let $E_k(y)$ be the expected discrepancy on an assignment with $k$ ones. The exchange of the values $0$ and $1$ gives $E_k=E_{n-k}$, so $\OPT_n=\min_y\max_{1\le k\le n/2}E_k(y)$. For $n\le3$ only $E_1$ occurs. For $n=4$ and $n=5$ both $E_1$ and $E_2$ occur, and for $p$ above a threshold ($\alpha_4\approx0.4245$ for $n=4$) the optimal algorithm is a mixture $tu+(1-t)v$ of two fixed algorithms $u$ and $v$ with a coefficient $t$ that depends on $p$.

On $[\alpha_4,\h]$ and on $[\h,1]$ the function $\OPT_4$ is rational and not a polynomial, whereas the expected discrepancy of a fixed algorithm on a fixed input assignment is a polynomial in $p$. Hence on these intervals no finite set of algorithms that do not depend on $p$ contains an optimal algorithm for every $p$. Each $E_k$ is convex on $[0,1]^m$ and linear on each cell of the subdivision of the cube by the hyperplanes $y_i=y_j$, $y_i+y_j=1$ and $y_i=\h$, and every vertex of this subdivision lies in $\{0,\h,1\}^m$ (Lemma~\ref{lem:vert}). Hence, for every $\lambda\in[0,1]$,
\[
\OPT_n\;\ge\;\min_{y\in[0,1]^m}\bigl(\lambda E_1+(1-\lambda)E_2\bigr)(y)\;=\;\min_{w\in\{0,\h,1\}^m}\bigl(\lambda E_1+(1-\lambda)E_2\bigr)(w),\qquad n=4,5
\]
(Proposition~\ref{prop:lower}). For $w\in\{0,\h,1\}^m$ the $E_k(w)$ are polynomials in $p$ with rational coefficients. When $\lambda$ is a rational function of $p$, a bound $\OPT_n\ge\Phi$ on an interval of $p$ follows from $3^m$ inequalities between rational functions on that interval, and these are decided by Sturm sequences. The matching upper bound comes from the mixture $tu+(1-t)v$ with $t$ chosen so that $tE_k(u)+(1-t)E_k(v)$ takes the same value for $k=1$ and $k=2$; by convexity this value bounds the cost of the mixture (Lemma~\ref{lem:mix}).

In a related model a broadcast reaches all other processes or none of them; consensus in that model is studied in \cite{FPR26}. Knight, Lehnherr and Rajsbaum \cite{KLR26} treat randomized versions of coordinated attack, a problem posed by Gray \cite{Gray}, and of approximate agreement and consensus with lost messages in probabilistic epistemic logic. We keep the model of \cite{FPR}, in which each message is lost independently.

An algorithm is symmetric if it is invariant under permutations of the processes and under the exchange of the values $0$ and $1$. A symmetric algorithm is described by the value $y(k_0,k_1)$ decided by a process with input $0$ that received $k_0$ zeros and $k_1\ge1$ ones. A process with input $0$ that received no one must decide $0$, and a process with input $1$ decides $1-y(k_1,k_0)$. We call the pairs $(k_0,k_1)$ views, and the vector of the values $y(k_0,k_1)$ over the views, in the order below, the decision vector of the algorithm. For $n=4$ the views are
\[
(0,1),\ (0,2),\ (0,3),\ (1,1),\ (1,2),\ (2,1),
\]
and for $n=5$ they are $(0,1),(0,2),(0,3),(0,4),(1,1),(1,2),(1,3),(2,1),(2,2),(3,1)$. For every $n$, FV decides $1$ in every view, AMP$(a)$ decides $a$ in every view, and $\mathrm{AMP}'$ is AMP$(\h)$ except that $y(0,n-1)=1$. For $n=4$ these and one more decision vector are
\begin{align*}
\mathrm{FV}&=(1,1,1,1,1,1), & \mathrm{AMP}(\h)&=(\h,\h,\h,\h,\h,\h),\\
\mathrm{AMP}'&=(\h,\h,1,\h,\h,\h), & v_4&=(\h,1,1,0,\h,0).
\end{align*}
The vector $v_4$ decides by majority among the values a process knows, its own included, and decides $\h$ on a tie.

\begin{theorem}\label{thm:four}
Let $\beta_4=1-2^{-1/4}\approx0.159104$ and let $\alpha_4\approx0.424520$ be the root in $(0,1)$ of $2p^5-3p^4+5p^3-11p^2+11p-3$. Then
\[
\OPT_4(p)=\begin{cases}
1-2q^4(1-q^2)^2, & 0\le p\le\beta_4,\\
1-(1-q^2)^2, & \beta_4\le p\le\alpha_4,\\
\Phi_1(p), & \alpha_4\le p\le\h,\\
\Phi_{1/2}(p), & \h\le p\le1,
\end{cases}
\]
where $\Phi_1$ and $\Phi_{1/2}$ are the rational functions of Table~\ref{tab:four}. The optimum is attained by FV on the first interval, by AMP$(\h)$ for $\beta_4\le p\le 2-\sqrt{10}/2\approx0.418861$ and by AMP$'$ for $0.418861\le p\le\alpha_4$. On the third interval it is attained by $(1-2t_1)\,\mathrm{AMP}'+2t_1v_4$ and on the fourth by $(1-2t_{1/2})\,\mathrm{AMP}(\h)+2t_{1/2}v_4$, with $t_1,t_{1/2}$ as in Table~\ref{tab:four}. In particular $\OPT_4(\h)=941/1936$, attained with $t_{1/2}(\h)=8/121$.
\end{theorem}

\begin{table}[t]
\centering
\small
\begin{tabular}{l}
\toprule
$D_1=12p^{10}-50p^9+84p^8-49p^7-63p^6+167p^5-168p^4+85p^3-12p^2-4p-3$\\
$\Phi_1=(p-1)(16p^{15}-88p^{14}+240p^{13}-444p^{12}+682p^{11}-967p^{10}+1222p^9-1234p^8$\\
\qquad$+909p^7-435p^6+106p^5-3p^4+9p^3-17p^2+7p+3)/D_1$\\
$t_1=-(2p^5-3p^4+5p^3-11p^2+11p-3)/(2pD_1)$\\
\midrule
$D_{1/2}=12p^9-38p^8+46p^7-3p^6-66p^5+101p^4-65p^3+17p^2+8p+3$\\
$\Phi_{1/2}=(p-1)(4p^{12}-54p^{11}+224p^{10}-501p^9+691p^8-616p^7+333p^6-81p^5$\\
\qquad$-14p^4+2p^3+8p^2-11p-3)/D_{1/2}$\\
$t_{1/2}=-(2p^2-8p+3)/(2pD_{1/2})$\\
\bottomrule
\end{tabular}
\caption{The formulas of Theorem~\ref{thm:four} on $[\alpha_4,\h]$ and on $[\h,1]$.}\label{tab:four}
\end{table}

Below $\beta_4$ the optimum is attained by FV, and from $\beta_4$ to $0.4189$ by AMP$(\h)$, as for two processes. Above $\alpha_4$ neither is optimal. The algorithm of Theorem~\ref{thm:four} moves the decisions in the views $(0,2)$, $(1,1)$ and $(2,1)$ away from $\h$ by $t_1(p)$ on the third interval and by $t_{1/2}(p)$ on the fourth; this displacement is $0$ at $\alpha_4$ and at most about $0.124$, and on the fourth interval the decision in the view $(0,3)$ moves with it. The subscript of $\Phi_1$ and $t_1$ is the decision of $\mathrm{AMP}'$ in the view $(0,3)$, and that of $\Phi_{1/2}$ and $t_{1/2}$ the decision of AMP$(\h)$ there. The formulas of the third and fourth intervals agree at $p=\h$, where both give $941/1936$. For $n=2$ and $n=3$ the optimum equals $q$ for $p\ge\h$; for $n=4$ it is below $q$ up to a point $p_0\approx0.640$ and above $q$ from there to $1$ (Table~\ref{tab:values}).

For $n=5$ the optimal algorithms are given in Table~\ref{tab:five}. There $\mathrm{FV}_1$ decides $1$ in the views with at most one zero and $0$ in the others.

\begin{theorem}\label{thm:five}
There are $0<\beta_5<\gamma_5<\delta_5<\h<\epsilon_5<\zeta_5<1$, with $\beta_5\approx0.109884$, $\gamma_5\approx0.410601$, $\delta_5\approx0.414398$, $\epsilon_5\approx0.565484$ and $\zeta_5\approx0.673303$, algebraic numbers whose minimal polynomials are given in Table~\ref{tab:five}, such that
\[
\OPT_5(p)=\begin{cases}
P_5(q), & 0\le p\le\beta_5,\\[2pt]
1-\h\bigl[(1-q^2)^3+(1-q^3)^2\bigr], & \beta_5\le p\le\gamma_5,
\end{cases}
\]
where
\[
P_5(q)=q\,(6-15q+16q^2-12q^4+33q^7-42q^8-3q^9+42q^{10}-25q^{11}-6q^{12}+9q^{13}-2q^{14}).
\]
On each of the intervals $[\gamma_5,\delta_5]$, $[\delta_5,\h]$, $[\h,\epsilon_5]$, $[\epsilon_5,\zeta_5]$ and $[\zeta_5,1]$, $\OPT_5$ is the rational function
\begin{equation}\label{eq:mix}
\Phi_{u,v}=\frac{a_ub_v-a_vb_u}{(a_u-b_u)-(a_v-b_v)},\qquad a_w=E_1(w),\ b_w=E_2(w),
\end{equation}
taken in lowest terms, for the pair $(u,v)$ of Table~\ref{tab:five}, where $E_k(w)$ is the expected discrepancy of the decision vector $w$ on an assignment with $k$ ones; for these pairs $2E_1$ and $2E_2$ are given in Table~\ref{tab:poly5}. The optimum is attained by $\mathrm{FV}_1$ on the first interval, by $\mathrm{AMP}'$ on the second, and on the others by $tu+(1-t)v$ with $t=(b_v-a_v)/((a_u-b_u)+(b_v-a_v))\in[0,1]$, again in lowest terms. In particular $\OPT_5(\h)=394955/851328\approx0.463928$, and $\OPT_5(1)=0$.
\end{theorem}

\begin{table}[t]
\centering
\small
\renewcommand{\arraystretch}{1.35}
\begin{tabular}{llcccccccccc}
\toprule
interval & & $(0,1)$&$(0,2)$&$(0,3)$&$(0,4)$&$(1,1)$&$(1,2)$&$(1,3)$&$(2,1)$&$(2,2)$&$(3,1)$\\
\midrule
$[0,\beta_5]$ & $\mathrm{FV}_1$ & 1&1&1&1&1&1&1&0&0&0\\
$[\beta_5,\gamma_5]$ & $\mathrm{AMP}'$ & $\h$&$\h$&$\h$&1&$\h$&$\h$&$\h$&$\h$&$\h$&$\h$\\
$[\gamma_5,\delta_5]$ & $u=\mathrm{AMP}'$ & $\h$&$\h$&$\h$&1&$\h$&$\h$&$\h$&$\h$&$\h$&$\h$\\
 & $v$ & $\h$&$\h$&1&1&$\h$&$\h$&$\h$&$\h$&$\h$&$\h$\\
$[\delta_5,\h]$ & $u=\mathrm{AMP}'$ & $\h$&$\h$&$\h$&1&$\h$&$\h$&$\h$&$\h$&$\h$&$\h$\\
 & $v$ & $\h$&1&1&1&0&$\h$&1&0&$\h$&0\\
$[\h,\epsilon_5]$ & $u=\mathrm{AMP}(\h)$ & \multicolumn{10}{c}{all $\h$}\\
 & $v$ & $\h$&1&1&1&0&$\h$&1&0&$\h$&0\\
$[\epsilon_5,\zeta_5]$ & $u=\mathrm{AMP}(\h)$ & \multicolumn{10}{c}{all $\h$}\\
 & $v$ & $\h$&1&1&1&0&$\h$&1&0&0&0\\
$[\zeta_5,1]$ & $u=\mathrm{AMP}(\h)$ & \multicolumn{10}{c}{all $\h$}\\
 & $v$ & $\h$&1&1&1&0&$\h$&$\h$&0&$\h$&0\\
\bottomrule
\end{tabular}

\medskip
\begin{tabular}{ll}
\toprule
$\beta_5$ & $4p^{12}-38p^{11}+142p^{10}-246p^9+114p^8+306p^7-606p^6+414p^5+66p^4$\\
 & \qquad$-352p^3+284p^2-109p+9$\\
$\gamma_5$ & $4p^6-6p^5+4p^4+2p^3-14p^2+15p-4$\\
$\delta_5$ & $39p^{21}-651p^{20}+5035p^{19}-23855p^{18}+77068p^{17}-181614p^{16}+336505p^{15}$\\
 & \qquad$-544861p^{14}+852940p^{13}-1299598p^{12}+1779201p^{11}-2029409p^{10}$\\
 & \qquad$+1851023p^{9}-1318201p^{8}+712758p^{7}-276514p^{6}+66052p^{5}$\\
 & \qquad$-3788p^{4}-2634p^{3}+306p^{2}+276p-72$\\
$\epsilon_5$ & $249p^{10}-1504p^9+4107p^8-6636p^7+6953p^6-4830p^5+2130p^4-468p^3-48p^2+60p-12$\\
$\zeta_5$ & $325p^9-1803p^8+4488p^7-6480p^6+5954p^5-3592p^4+1389p^3-291p^2+8p+6$\\
\bottomrule
\end{tabular}
\caption{Optimal one-round algorithms for $n=5$, as the decision of a process with input $0$ in each view $(k_0,k_1)$, and the minimal polynomials of the thresholds.}\label{tab:five}
\end{table}

\begin{table}[t]
\centering
\small
\begin{tabular}{cccc}
\toprule
$p$ & $\OPT_4$ & $\OPT_5$ & $q$\\
\midrule
0.05 & 0.984514 & 0.984358 & 0.95\\
0.1 & 0.952630 & 0.957158 & 0.9\\
0.2 & 0.870400 & 0.857600 & 0.8\\
0.3 & 0.739900 & 0.717850 & 0.7\\
0.4 & 0.590400 & 0.561600 & 0.6\\
0.5 & 0.486054 & 0.463928 & 0.5\\
0.6 & 0.396044 & 0.378065 & 0.4\\
0.7 & 0.305296 & 0.294896 & 0.3\\
0.8 & 0.211238 & 0.210054 & 0.2\\
0.9 & 0.111271 & 0.116181 & 0.1\\
\bottomrule
\end{tabular}
\caption{Values of $\OPT_4$ and $\OPT_5$ at selected $p$, rounded to six decimals, with $q=1-p$, the optimum for two and for three processes when $p\ge\h$, for comparison.}\label{tab:values}
\end{table}

For five processes and $p\le\beta_5$ the optimal vector $\mathrm{FV}_1$ decides $1$ only in the views with at most one zero. Between $\beta_5$ and $\gamma_5$ the optimum equals the expected discrepancy of AMP$(\h)$ on assignments with two ones, and $\mathrm{AMP}'$ keeps its expected discrepancy on assignments with one $1$ below that value. Above $\gamma_5$ the algorithm of Theorem~\ref{thm:five} mixes two decision vectors with entries in $\{0,\h,1\}$, as above $\alpha_4$ for four processes, and the decisions depend on $p$.

The optimal one-round algorithms for $n\ge6$ processes, and the optimal algorithms with more than one round for $n\ge3$, are not known.

\section{Reduction}

Fix $n$ and $p$. Up to a permutation of the processes, an input assignment is determined by its number $k$ of ones, and by validity only $1\le k\le n-1$ matter.

\begin{lemma}\label{lem:sym}
For every one-round algorithm there is a symmetric one, given by a function $y(k_0,k_1)$ as in Section~\ref{sec:intro}, whose expected discrepancy on every input assignment is at most the cost of the original. For a symmetric algorithm the expected discrepancy $E_k$ on an assignment with $k$ ones satisfies $E_k=E_{n-k}$.
\end{lemma}

\begin{proof}
The discrepancy $\max_id_i-\min_id_i$ is a convex function of the vector $(d_1,\dots,d_n)$ of decisions. Average the decisions of the algorithm over all permutations of the processes and over the exchange of the values $0$ and $1$, which replaces every decision $y$ by $1-y$ and every received value $x$ by $1-x$. The average is again a deterministic one-round algorithm and satisfies validity. By convexity its expected discrepancy on each assignment is at most the average of the original's over the orbit of that assignment, hence at most the original's cost. It is invariant under permutations of the processes, so a decision depends only on the input and on how many zeros and ones were received. The exchange gives $y(1,k_0,k_1)=1-y(0,k_1,k_0)$ and $E_k=E_{n-k}$. A process with input $0$ that received only zeros, or nothing, must decide $0$ by validity, since its view is consistent with all inputs being $0$.
\end{proof}

So $\OPT_n=\min_y\max_{1\le k\le n/2}E_k(y)$ over $y\in[0,1]^m$, as stated in Section~\ref{sec:intro}, and for $n=4$ and $n=5$ the maximum is that of $E_1$ and $E_2$.

\begin{lemma}\label{lem:vert}
Each $E_k$ is a convex piecewise linear function of $y$ on $[0,1]^m$, linear on the cells of the subdivision of the cube by the hyperplanes $y_i=y_j$, $y_i+y_j=1$ and $y_i=\h$. Every vertex of this subdivision lies in $\{0,\h,1\}^m$. Consequently, for every $\lambda\in[0,1]$ the minimum of $\lambda E_1+(1-\lambda)E_2$ over the cube is attained in $\{0,\h,1\}^m$.
\end{lemma}

\begin{proof}
In a fixed execution each decision is $0$, $1$, $y_i$ or $1-y_i$ for some view $i$. Once the order of the numbers $0,1,y_i,1-y_i$ is fixed, the maximum and the minimum are fixed affine functions of $y$, and this order is constant on each cell of the subdivision. Taking expectations, each $E_k$ is linear on each cell. It is convex as an expectation of maxima minus minima of affine functions.

A vertex is cut out by $m$ independent equations, each of the form $y_i=c$ with $c\in\{0,\h,1\}$ or $y_i=\pm y_j+c$ with $c\in\{0,1\}$. Join two coordinates when an equation of the second form relates them. On a connected class, every coordinate has the form $y_i=s_iz+c_i$ with $s_i=\pm1$, $c_i\in\mathbb{Z}$ and one free parameter $z$. The class is pinned either by an equation of the first form, which gives $z\in\h\mathbb{Z}$, or by a cycle whose equation is not implied by the others, which reads $2z\in\mathbb{Z}$. In both cases every coordinate of the class lies in $\h\mathbb{Z}\cap[0,1]=\{0,\h,1\}$. A convex function that is linear on each cell of a polyhedral subdivision of the cube attains its minimum at a vertex of the subdivision.
\end{proof}

\begin{proposition}\label{prop:lower}
Suppose that on an interval $I$ of values of $p$ there are $\lambda(p)\in[0,1]$ and $\Phi(p)$ with $\lambda E_1(w)+(1-\lambda)E_2(w)\ge\Phi$ for every $w\in\{0,\h,1\}^m$. Then $\OPT_n\ge\Phi$ on $I$.
\end{proposition}

\begin{proof}
For every $y$, $\max(E_1,E_2)(y)\ge\lambda E_1(y)+(1-\lambda)E_2(y)$, and by Lemma~\ref{lem:vert} the right side is at least its minimum over $\{0,\h,1\}^m$.
\end{proof}

By Sion's minimax theorem \cite{Sion} and Lemma~\ref{lem:vert} the best bound of this kind is exact for $n=4,5$: $\OPT_n=\max_\lambda\min_w(\lambda E_1+(1-\lambda)E_2)(w)$.

\begin{lemma}\label{lem:mix}
Let $u,v$ be decision vectors such that $a_u-b_u$ and $a_v-b_v$ are not both zero and do not have the same strict sign, where $a_w=E_1(w)$ and $b_w=E_2(w)$. Put $t=(b_v-a_v)/((a_u-b_u)+(b_v-a_v))$. Then $y=tu+(1-t)v$ has $\max(E_1,E_2)(y)\le\Phi_{u,v}$, with $\Phi_{u,v}$ as in \eqref{eq:mix}. Moreover, with $\lambda=(b_v-b_u)/((a_u-a_v)+(b_v-b_u))$ we have $\lambda a_u+(1-\lambda)b_u=\lambda a_v+(1-\lambda)b_v=\Phi_{u,v}$.
\end{lemma}

\begin{proof}
By hypothesis $a_u-b_u$ and $b_v-a_v$ have the same weak sign and are not both zero, so $t\ge0$ and $1-t\ge0$. By convexity $E_k(y)\le tE_k(u)+(1-t)E_k(v)$ for $k=1,2$. Write $\Delta=(a_u-b_u)+(b_v-a_v)$, so that $t=(b_v-a_v)/\Delta$ and $1-t=(a_u-b_u)/\Delta$. Then
\begin{align*}
ta_u+(1-t)a_v&=\frac{(b_v-a_v)a_u+(a_u-b_u)a_v}{\Delta}=\frac{a_ub_v-a_vb_u}{\Delta},\\
tb_u+(1-t)b_v&=\frac{(b_v-a_v)b_u+(a_u-b_u)b_v}{\Delta}=\frac{a_ub_v-a_vb_u}{\Delta},
\end{align*}
and $\Delta=(a_u-b_u)-(a_v-b_v)$, so both equal $\Phi_{u,v}$. For the second statement put $\Lambda=(a_u-a_v)+(b_v-b_u)$, so that $\lambda=(b_v-b_u)/\Lambda$ and $1-\lambda=(a_u-a_v)/\Lambda$. Then $\lambda a_u+(1-\lambda)b_u=(a_ub_v-a_vb_u)/\Lambda=\lambda a_v+(1-\lambda)b_v$, and $\Lambda=\Delta$.
\end{proof}

\section{Proofs of Theorems \ref{thm:four} and \ref{thm:five}}

For a decision vector $w\in\{0,\h,1\}^m$ the functions $E_k(w)$ are explicit polynomials. Let $L=n(n-1)$ be the number of messages in a round. Then $2E_k(w)=\sum_{j=0}^{L}M_k[w,j]\,p^jq^{L-j}$, where the integer $M_k[w,j]$ is the sum of the doubled discrepancies of the executions in which exactly $j$ messages are delivered. The matrices $M_k$ are computed by enumerating, for each receiving process, the numbers of zeros and ones it receives, combining the receivers, and grouping executions by the set of decisions they produce. All the inequalities below are decided in exact rational and algebraic arithmetic in Python, with SymPy~1.14 \cite{SymPy}.

\emph{Lower bounds.} On each interval of the two theorems we apply Proposition~\ref{prop:lower} with a function $\lambda(p)$ and the value $\Phi(p)$ stated there. On the first two intervals of each theorem $\lambda=0$. On the others $\lambda$ is the rational function of Lemma~\ref{lem:mix} for the pair $(u,v)$ that attains the optimum. For $n=4$, with
\[
R=8p^8-32p^7+58p^6-54p^5+18p^4+16p^3-15p^2-p+8,
\]
it is $\lambda=2p(p-1)R/D_1$ on $[\alpha_4,\h]$ and $\lambda=2pR/D_{1/2}$ on $[\h,1]$, with $D_1$ and $D_{1/2}$ as in Table~\ref{tab:four}. On each interval $0\le\lambda\le1$, and for each $w\in\{0,\h,1\}^m$ the rational function $\lambda E_1(w)+(1-\lambda)E_2(w)-\Phi$, taken in lowest terms, is nonnegative; these are $3^6=729$ inequalities on each interval for $n=4$ and $3^{10}=59049$ for $n=5$. The denominator has no zero on the interval, so each inequality is the nonnegativity of a polynomial with rational coefficients on an interval whose endpoints are rational or real algebraic. Sturm sequences count its real roots in a rational interval enclosing $I$. The roots inside $I$ are isolated as algebraic numbers and compared with the endpoints through their minimal polynomials and isolating intervals. Finally the sign is read at a rational point of each subinterval between consecutive roots. Decision vectors with the same pair of polynomials give the same condition.

\emph{Upper bounds.} At $p=1$ every message is delivered, AMP$(\h)$ has discrepancy $0$, and $\OPT_n(1)=0$; below we take $p<1$. For $n=4$, $E_2(\mathrm{FV})=1-2q^4(1-q^2)^2$ and $E_2(\mathrm{AMP}(\h))=E_2(\mathrm{AMP}')=1-(1-q^2)^2$, and
\begin{align*}
E_2(\mathrm{FV})-E_1(\mathrm{FV})&=pq^3(2p^4-10p^3+18p^2-11p+3),\\
E_2(\mathrm{AMP}(\h))-E_1(\mathrm{AMP}(\h))&=\h pq(2p^2-8p+3),\\
E_2(\mathrm{AMP}')-E_1(\mathrm{AMP}')&=-\h p(2p^5-3p^4+5p^3-11p^2+11p-3).
\end{align*}
The quartic has no real root, $2p^2-8p+3\ge0$ for $p\le2-\sqrt{10}/2$, and the quintic is negative on $[0,\alpha_4)$. Hence $E_1\le E_2$ for FV on $[0,1]$, for AMP$(\h)$ on $[0,2-\sqrt{10}/2]$ and for $\mathrm{AMP}'$ on $[0,\alpha_4]$. On the third and fourth intervals Lemma~\ref{lem:mix} applies to the pairs $(\mathrm{AMP}',v_4)$ and $(\mathrm{AMP}(\h),v_4)$. The coefficient it puts on $v_4$ is $2t_1$, respectively $2t_{1/2}$, and $\Phi_{u,v}$ equals $\Phi_1$, respectively $\Phi_{1/2}$. At $p=\h$ both pairs give $\lambda=94/121$. For $n=5$ the vectors $\mathrm{FV}_1$ and $\mathrm{AMP}'$ have $E_2$ equal to the stated optimum and $E_1$ at most it on their intervals. On the other intervals Lemma~\ref{lem:mix} applies, since the sign condition on $a_u-b_u$ and $a_v-b_v$ holds there. At $p=\h$ both adjacent pairs give $394955/851328$, with $\lambda=49103/79812$.

By \cite[proof of Theorem~14]{FPR}, for every $n$,
\[
E_k(\mathrm{AMP}(\h))=1-\h\bigl[(1-q^k)^{n-k}+(1-q^{n-k})^k\bigr],
\]
which gives the second line of both theorems.

\section*{Funding}
The research presented in this paper was supported by the European Research Council (ERC) under the European Union's Horizon Europe research and innovation programme (grant agreement No.~101041711), by the Simons Foundation as part of the Collaboration on the Mathematical and Scientific Foundations of Deep Learning, by Heights Labs, by Convex Nexus Capital, by the Israel Science Foundation (grant number 2258/19), and by the Israel Science Foundation (ISF Grant 4101/25).

\section*{Acknowledgments}
This work was done by Claude (Anthropic), connected to Math Brain~\cite{MathBrainRepo} and working inside the Brain framework~\cite{BrainFramework}, a shared, persistent memory in which many Claude and Codex agents record claims, proofs, failed attempts and methods, retrieve one another's results, coordinate who works on what, and audit each other's work. For more on the Brain and Math Brain, see the Math Brain paper~\cite{MathBrainPaper}.

\section*{Code and data availability}
The code and data are at \url{https://github.com/yspennstate/one-round-agreement-four-five}.
\section*{Declaration of generative AI and AI-assisted technologies}
This paper is part of a broader research project on AI systems proving mathematical results. The computations, proofs and text were produced by Claude (Anthropic). The model, the algorithms FV and AMP, and the solutions for two and three processes are due to Fraigniaud, Patt-Shamir and Rajsbaum \cite{FPR}, who asked for the optimal algorithms for every $n>2$; the results here extend theirs to four and five processes in one round. The author directed the work and is responsible for the content of the paper and for decisions concerning its dissemination.

\section*{Declaration of competing interest}
The author holds equity in Heights Labs and in Convex Nexus Capital and has a management role at Convex Nexus Capital. These entities are included in the funding declaration above.

\appendix
\section{The polynomials for five processes}

Table~\ref{tab:poly5} gives $2E_1$ and $2E_2$ for the decision vectors of Table~\ref{tab:five} on the last five intervals of Theorem~\ref{thm:five}, with $q=1-p$. The vectors $v_{\gamma}$, $v_{\delta}$, $v_{\epsilon}$ and $v_{\zeta}$ are the rows labelled $v$ on $[\gamma_5,\delta_5]$, on $[\delta_5,\h]$ and $[\h,\epsilon_5]$, on $[\epsilon_5,\zeta_5]$, and on $[\zeta_5,1]$. The last three share
\begin{align*}
2E_1&=q^3S,\\
S&=13p^{13}-113p^{12}+432p^{11}-948p^{10}+1324p^9-1260p^8+882p^7\\
&\qquad-482p^6+201p^5-69p^4+24p^3+2p+2.
\end{align*}
With these, $\Phi_{u,v}$ of \eqref{eq:mix} gives $\OPT_5$ at every $p\ge\gamma_5$; at $p=\h$ the pairs $(\mathrm{AMP}',v_\delta)$ and $(\mathrm{AMP}(\h),v_\delta)$ both give $394955/851328$.

\begin{table}[ht]
\centering
\small
\begin{tabular}{l}
\toprule
$\mathrm{AMP}(\h)$:\quad $2E_1=2q(2p^2-p+1)$,\quad $2E_2=q^2(2+4p-3p^2)$\\
$\mathrm{AMP}'$:\quad $2E_1=4p^7-6p^6+4p^5-p^4-4p^3+6p^2-4p+2$,\quad $2E_2=q^2(2+4p-3p^2)$\\
\midrule
$v_\gamma$:\quad $2E_1=-12p^7+34p^6-36p^5+19p^4-8p^3+6p^2-4p+2$\\
\qquad $2E_2=q(3p^{11}-15p^{10}+25p^9-9p^8-17p^7+15p^6+3p^3-7p^2+2p+2)$\\
\midrule
$v_\delta$:\quad $2E_1=q^3S$\\
\qquad $2E_2=-189p^{16}+1804p^{15}-7719p^{14}+19722p^{13}-33720p^{12}+40764p^{11}-35649p^{10}$\\
\qquad\qquad $+22434p^9-9666p^8+2382p^7-33p^6-150p^5+28p^3-9p^2+2$\\
\midrule
$v_\epsilon$:\quad $2E_1=q^3S$\\
\qquad $2E_2=q(438p^{15}-2870p^{14}+8956p^{13}-17402p^{12}+23271p^{11}-22323p^{10}+15456p^9$\\
\qquad\qquad $-7446p^8+2172p^7-150p^6-129p^5+21p^4+21p^3-7p^2+2p+2)$\\
\midrule
$v_\zeta$:\quad $2E_1=q^3S$\\
\qquad $2E_2=q(113p^{15}-1067p^{14}+4468p^{13}-10922p^{12}+17317p^{11}-18731p^{10}+14067p^9$\\
\qquad\qquad $-7155p^8+2164p^7-156p^6-129p^5+21p^4+21p^3-7p^2+2p+2)$\\
\bottomrule
\end{tabular}
\caption{The polynomials $2E_1$ and $2E_2$ for the decision vectors of Table~\ref{tab:five} above $\gamma_5$.}\label{tab:poly5}
\end{table}

\begin{thebibliography}{9}
\bibitem{FPR} P.~Fraigniaud, B.~Patt-Shamir and S.~Rajsbaum, Coordination through stochastic channels, in \emph{39th International Symposium on Distributed Computing (DISC 2025)}, LIPIcs 356 (2025), 32:1--32:19, doi:10.4230/LIPIcs.DISC.2025.32.
\bibitem{FPR26} P.~Fraigniaud, B.~Patt-Shamir and S.~Rajsbaum, Consensus with stochastic broadcast, in \emph{40th International Symposium on Distributed Computing (DISC 2026)}, LIPIcs 397 (2026), 18:1--18:17, doi:10.4230/LIPIcs.DISC.2026.18; arXiv:2608.27336.
\bibitem{DLPSW} D.~Dolev, N.~A.~Lynch, S.~S.~Pinter, E.~W.~Stark and W.~E.~Weihl, Reaching approximate agreement in the presence of faults, \emph{J. ACM} 33 (1986), 499--516, doi:10.1145/5925.5931.
\bibitem{Gray} J.~N.~Gray, Notes on data base operating systems, in \emph{Operating Systems: An Advanced Course}, Lecture Notes in Computer Science 60, Springer, 1978, 393--481, doi:10.1007/3-540-08755-9\_9.
\bibitem{KLR26} S.~Knight, D.~Lehnherr and S.~Rajsbaum, An epistemic analysis of random coordinated attack, preprint, 2026, arXiv:2606.17860.
\bibitem{Sion} M.~Sion, On general minimax theorems, \emph{Pacific J. Math.} 8 (1958), 171--176.
\bibitem{SymPy} A.~Meurer, C.~P.~Smith, M.~Paprocki, O.~\v{C}ert\'{\i}k, S.~B.~Kirpichev, M.~Rocklin, A.~Kumar, S.~Ivanov, J.~K.~Moore, S.~Singh, T.~Rathnayake, S.~Vig, B.~E.~Granger, R.~P.~Muller, F.~Bonazzi, H.~Gupta, S.~Vats, F.~Johansson, F.~Pedregosa, M.~J.~Curry, A.~R.~Terrel, \v{S}.~Rou\v{c}ka, A.~Saboo, I.~Fernando, S.~Kulal, R.~Cimrman and A.~Scopatz, SymPy: symbolic computing in Python, \emph{PeerJ Comput. Sci.} 3 (2017), e103, doi:10.7717/peerj-cs.103.
\bibitem{MathBrainRepo} Y.~Shmalo, \emph{Math Brain: public snapshot, code and tests}, 2026, \url{https://github.com/yspennstate/math-brain-public}.
\bibitem{MathBrainPaper} Y.~Shmalo, \emph{Math Brain: shared, auditable memory for multi-agent mathematical research}, October 2026, doi:10.5281/zenodo.24150837, \url{https://github.com/yspennstate/math-brain-public/tree/main/papers/math-brain-2026-10-24}.
\bibitem{BrainFramework} Y.~Shmalo, \emph{Brain Framework}, public software repository, 2026, \url{https://github.com/yspennstate/brain-framework-public}.
\end{thebibliography}
\end{document}
